% GENERATED FILE. Edit canonical lessons, then run npm run generate:books.
% canonical-source-hash: fnv1a64:893766cd4bbaf5b1
% canonical-lessons: MW-C235-guvaro, MW-C235-charo, MW-C235-gobar, MW-C235-palno, MW-C235-bhuso

\chapter{The Cattle Yard}
\label{ch:mw-the-cattle-yard}
\hlchaptermodality{235}

\begin{chapteropening}
\textbf{By the end of this chapter:} \emph{I can name the cattle yard, fodder and dung, and say I keep goats.}

Complete the last lesson of chapter 235: I can name the cattle yard, fodder and dung, and say I keep goats.
\end{chapteropening}

\section[\texorpdfstring{\mw{गुवाड़ो} (guvāṛo)}{गुवाड़ो (guvāṛo)}]{\mw{गुवाड़ो} (guvāṛo) — a cattle yard; a cowshed}

\label{lesson:MW-C235-guvaro}



\begin{quote}
Before the new one: say the Marwadi for to drive (animals or a cart), then the Marwadi for to look after.
\end{quote}

\subsection*{You'll want to know: \mw{गुवाड़ो}}
\textbf{\mw{गुवाड़ो}} — \emph{guvāṛo} — \textquotedblleft{}a cattle yard; a cowshed\textquotedblright{}.

\begin{quote}
\mw{गुवाड़ो} \mw{घर} \mw{कनै} \mw{है।}
\end{quote}

\begin{quote}
\emph{guvāṛo ghar kanai hai.}
\end{quote}

\begin{quote}
The cattle yard is near the house.
\end{quote}

A masculine noun.

\begin{cousinweb}
\textbf{\mw{गुवाड़ो}} (\emph{guvāṛo}) is inherited from Sanskrit \emph{govāṭa}, \textquotedblleft{}a fold for cattle\textquotedblright{}, built from \emph{go}, \textquotedblleft{}cow\textquotedblright{}, and \emph{vāṭa}, \textquotedblleft{}an enclosure\textquotedblright{}; CDIAL lists Marwari \emph{guvāṛo}, \textquotedblleft{}a cowshed\textquotedblright{} (CDIAL 4327).
\end{cousinweb}

\begin{grammarlens}[title={What you've built}]
One more word for this chapter.
\end{grammarlens}

\subsection*{Your turn}
\begin{itemize}
\raggedright
  \item \emph{Say it:} \emph{guvāṛo}
  \item \emph{Say it:} \emph{guvāṛo}, once more
  \item \emph{Say it:} \emph{guvāṛo ghar kanai hai.} — The cattle yard is near the house.
  \item \emph{Recall:} say the Marwadi for to drive (animals or a cart), then the Marwadi for to look after, then say \emph{guvāṛo} again
\end{itemize}

\begin{tcolorbox}[breakable,colback=teal!4,colframe=teal!35!black,title={Before you move on}]
What does \mw{गुवाड़ो} mean? (\textquotedblleft{}A cattle yard; a cowshed\textquotedblright{}.) Say it once more.
\end{tcolorbox}
\section[\texorpdfstring{\mw{चारो} (chāro)}{चारो (chāro)}]{\mw{चारो} (chāro) — fodder}

\label{lesson:MW-C235-charo}



\begin{quote}
Before the new one: say the Marwadi for to look after, then the Marwadi for a cattle yard.
\end{quote}

\subsection*{You'll want to know: \mw{चारो}}
\textbf{\mw{चारो}} — \emph{chāro} — \textquotedblleft{}fodder\textquotedblright{}.

\begin{quote}
\mw{भैंस} \mw{चारो} \mw{खावै} \mw{है।}
\end{quote}

\begin{quote}
\emph{bhaĩs chāro khāvai hai.}
\end{quote}

\begin{quote}
The buffalo eats fodder.
\end{quote}

A masculine noun.

\begin{cousinweb}
\textbf{\mw{चारो}} (\emph{chāro}) is inherited from Sanskrit \emph{cāra}, \textquotedblleft{}motion\textquotedblright{}, which CDIAL also takes in the sense of grazing and fodder; CDIAL lists a Gujarati word for fodder under it. CDIAL gives no Marwari form in this entry (CDIAL 4755).
\end{cousinweb}

\begin{grammarlens}[title={What you've built}]
One more word for this chapter.
\end{grammarlens}

\subsection*{Your turn}
\begin{itemize}
\raggedright
  \item \emph{Say it:} \emph{chāro}
  \item \emph{Say it:} \emph{chāro}, once more
  \item \emph{Say it:} \emph{bhaĩs chāro khāvai hai.} — The buffalo eats fodder.
  \item \emph{Recall:} say the Marwadi for to look after, then the Marwadi for a cattle yard, then say \emph{chāro} again
\end{itemize}

\begin{tcolorbox}[breakable,colback=teal!4,colframe=teal!35!black,title={Before you move on}]
What does \mw{चारो} mean? (\textquotedblleft{}Fodder\textquotedblright{}.) Say it once more.
\end{tcolorbox}
\section[\texorpdfstring{\mw{गोबर} (gobar)}{गोबर (gobar)}]{\mw{गोबर} (gobar) — cow dung}

\label{lesson:MW-C235-gobar}



\begin{quote}
Before the new one: say the Marwadi for a cattle yard, then the Marwadi for fodder.
\end{quote}

\subsection*{You'll want to know: \mw{गोबर}}
\textbf{\mw{गोबर}} — \emph{gobar} — \textquotedblleft{}cow dung\textquotedblright{}.

\begin{quote}
\mw{खेत} \mw{मांय} \mw{गोबर} \mw{है।}
\end{quote}

\begin{quote}
\emph{khet mā̃y gobar hai.}
\end{quote}

\begin{quote}
There's cow dung on the field.
\end{quote}

A masculine noun.

\begin{cousinweb}
\textbf{\mw{गोबर}} (\emph{gobar}) is inherited from Sanskrit \emph{govara}, \textquotedblleft{}powdered cow dung\textquotedblright{}; CDIAL lists Hindi \emph{gobar}, \textquotedblleft{}cow dung\textquotedblright{}. CDIAL gives no Marwari form in this entry (CDIAL 4316).
\end{cousinweb}

\begin{grammarlens}[title={What you've built}]
One more word for this chapter.
\end{grammarlens}

\subsection*{Your turn}
\begin{itemize}
\raggedright
  \item \emph{Say it:} \emph{gobar}
  \item \emph{Say it:} \emph{gobar}, once more
  \item \emph{Say it:} \emph{khet mā̃y gobar hai.} — There's cow dung on the field.
  \item \emph{Recall:} say the Marwadi for a cattle yard, then the Marwadi for fodder, then say \emph{gobar} again
\end{itemize}

\begin{tcolorbox}[breakable,colback=teal!4,colframe=teal!35!black,title={Before you move on}]
What does \mw{गोबर} mean? (\textquotedblleft{}Cow dung\textquotedblright{}.) Say it once more.
\end{tcolorbox}
\section[\texorpdfstring{\mw{पाळणो} (pāḷṇo)}{पाळणो (pāḷṇo)}]{\mw{पाळणो} (pāḷṇo) — to rear (animals); to keep}

\label{lesson:MW-C235-palno}



\begin{quote}
Before the new one: say the Marwadi for fodder, then the Marwadi for cow dung.
\end{quote}

\subsection*{You'll want to know: \mw{पाळणो}}
\textbf{\mw{पाळणो}} — \emph{pāḷṇo} — \textquotedblleft{}to rear (animals); to keep\textquotedblright{}.

\begin{quote}
\mw{हूं} \mw{बकरी} \mw{पाळूं} \mw{हूं।}
\end{quote}

\begin{quote}
\emph{hū̃ bakrī pāḷū̃ hū̃.}
\end{quote}

\begin{quote}
I keep goats.
\end{quote}

Past (I …): \mw{म्हैं} \mw{बकरी} \mw{पाळी}, \emph{mhaĩ bakrī pāḷī}.

\begin{cousinweb}
\textbf{\mw{पाळणो}} (\emph{pāḷṇo}) is inherited from Sanskrit \emph{pālayati}, \textquotedblleft{}guards\textquotedblright{}; CDIAL lists the Old Marwari verb \emph{pālaï}, \textquotedblleft{}to rear, to keep\textquotedblright{} (CDIAL 8129).
\end{cousinweb}

\begin{grammarlens}[title={What you've built}]
One more word for this chapter.
\end{grammarlens}

\subsection*{Your turn}
\begin{itemize}
\raggedright
  \item \emph{Say it:} \emph{pāḷṇo}
  \item \emph{Say it:} \emph{pāḷṇo}, once more
  \item \emph{Say it:} \emph{hū̃ bakrī pāḷū̃ hū̃.} — I keep goats.
  \item \emph{Recall:} say the Marwadi for fodder, then the Marwadi for cow dung, then say \emph{pāḷṇo} again
\end{itemize}

\begin{tcolorbox}[breakable,colback=teal!4,colframe=teal!35!black,title={Before you move on}]
What does \mw{पाळणो} mean? (\textquotedblleft{}To rear (animals); to keep\textquotedblright{}.) Say it once more.
\end{tcolorbox}
\section[\texorpdfstring{\mw{भूसो} (bhūso)}{भूसो (bhūso)}]{\mw{भूसो} (bhūso) — chaff; straw fodder}

\label{lesson:MW-C235-bhuso}



\begin{quote}
Before the new one: say the Marwadi for cow dung, then the Marwadi for to rear (animals).
\end{quote}

\subsection*{You'll want to know: \mw{भूसो}}
\textbf{\mw{भूसो}} — \emph{bhūso} — \textquotedblleft{}chaff; straw fodder\textquotedblright{}.

\begin{quote}
\mw{ऊंट} \mw{भूसो} \mw{खावै} \mw{है।}
\end{quote}

\begin{quote}
\emph{ū̃ṭ bhūso khāvai hai.}
\end{quote}

\begin{quote}
The camel eats chaff.
\end{quote}

A masculine noun.

\begin{cousinweb}
\textbf{\mw{भूसो}} (\emph{bhūso}) is inherited from Sanskrit \emph{busa}, \textquotedblleft{}chaff\textquotedblright{}; CDIAL lists Hindi \emph{bhūsā}, \textquotedblleft{}chaff, straw\textquotedblright{}, and a Gujarati cousin. CDIAL gives no Marwari form in this entry (CDIAL 9293).
\end{cousinweb}

\begin{grammarlens}[title={What you've built}]
One more word for this chapter.
\end{grammarlens}

\subsection*{Your turn}
\begin{itemize}
\raggedright
  \item \emph{Say it:} \emph{bhūso}
  \item \emph{Say it:} \emph{bhūso}, once more
  \item \emph{Say it:} \emph{ū̃ṭ bhūso khāvai hai.} — The camel eats chaff.
  \item \emph{Recall:} say the Marwadi for cow dung, then the Marwadi for to rear (animals), then say \emph{bhūso} again
\end{itemize}

\begin{tcolorbox}[breakable,colback=teal!4,colframe=teal!35!black,title={Before you move on}]
What does \mw{भूसो} mean? (\textquotedblleft{}Chaff; straw fodder\textquotedblright{}.) Say it once more.
\end{tcolorbox}
